% GENERATED FILE. Edit canonical lessons, then run npm run generate:books.
% canonical-source-hash: fnv1a64:a78e618b2e966e97
% canonical-lessons: SA-C278-mama-matena, SA-C278-yatah-karanat, SA-C278-manyate, SA-C278-drdhah, SA-S227-letter-ddha, SA-R278-opinion-and-reason

\chapter{What I Think, and Why}
\label{ch:sa-what-i-think-and-why}
\hlchaptermodality{278}

\begin{chapteropening}
\textbf{By the end of this chapter:} \emph{I can give an opinion (mama matena …, ahaṃ manye yat …) and a reason for it (yataḥ …, varṣākāraṇāt …).}

Complete SA-R278-opinion-and-reason: say what you think and give a reason for it.
\end{chapteropening}

\section[\texorpdfstring{\sk{मम} \sk{मतेन} (mama matena)}{मम मतेन (mama matena)}]{\sk{मम} \sk{मतेन} (mama matena) — in my opinion}

\label{lesson:SA-C278-mama-matena}



\begin{quote}
Before the new one: say \textquotedblleft{}an opinion\textquotedblright{}. (\emph{matam}.) And \textquotedblleft{}my\textquotedblright{}. (\emph{mama}.) This chapter gives your opinion, and a reason for it.
\end{quote}

\subsection*{You'll want to know: \sk{मम} \sk{मतेन}}
\begin{quote}
\textbf{\sk{मम} \sk{मतेन} \sk{एतत्} \sk{सुन्दरम्।}}
\end{quote}

\begin{quote}
\emph{mama matena etat sundaram.} — \textquotedblleft{}In my opinion this is beautiful.\textquotedblright{}
\end{quote}

\textbf{\sk{मतेन}} (\emph{matena}) is \textbf{\sk{मतम्}} (\emph{matam}), \textquotedblleft{}an opinion\textquotedblright{}, with \emph{-ena}, \textquotedblleft{}by\textquotedblright{}: \textquotedblleft{}by my opinion\textquotedblright{}, that is, \textquotedblleft{}in my opinion\textquotedblright{}.

\begin{quote}
\textbf{\sk{मम} \sk{मतेन} \sk{सः} \sk{सत्यं} \sk{वदति।}}
\end{quote}

\begin{quote}
\emph{mama matena saḥ satyaṃ vadati.} — \textquotedblleft{}In my opinion he is telling the truth.\textquotedblright{}
\end{quote}

\begin{quote}
\textbf{\sk{तव} \sk{मतेन} \sk{किम्}?}
\end{quote}

\begin{quote}
\emph{tava matena kim?} — \textquotedblleft{}What is your opinion?\textquotedblright{}
\end{quote}

\subsection*{You'll want to know: \sk{अहं} \sk{मन्ये}}
\textbf{\sk{मन्ये}} (\emph{manye}) is \textquotedblleft{}I think\textquotedblright{}. It ends in \emph{-e}, the \textquotedblleft{}I\textquotedblright{} ending of verbs like \textbf{\sk{रोचते}} (\emph{rocate}). \textbf{\sk{यत्}} (\emph{yat}), \textquotedblleft{}that\textquotedblright{}, opens what you think:

\begin{quote}
\textbf{\sk{अहं} \sk{मन्ये} \sk{यत्} \sk{वर्षा} \sk{भविष्यति।}}
\end{quote}

\begin{quote}
\emph{ahaṃ manye yat varṣā bhaviṣyati.} — \textquotedblleft{}I think that it will rain.\textquotedblright{}
\end{quote}

Or say the thought first and close it with \textbf{\sk{इति}} (\emph{iti}), as you close what someone said:

\begin{quote}
\textbf{\sk{वर्षा} \sk{भविष्यति} \sk{इति} \sk{मन्ये।}}
\end{quote}

\begin{quote}
\emph{varṣā bhaviṣyati iti manye.} — \textquotedblleft{}It will rain, I think.\textquotedblright{}
\end{quote}

\subsection*{Your turn}
Say these aloud:

\begin{itemize}
\raggedright
  \item \textquotedblleft{}I think that he will go\textquotedblright{} — \emph{ahaṃ manye yat saḥ gamiṣyati}
  \item \textquotedblleft{}I think that she will go\textquotedblright{} — \emph{ahaṃ manye yat sā gamiṣyati}
  \item \textquotedblleft{}I think that they will go\textquotedblright{} — \emph{ahaṃ manye yat te gamiṣyanti}
\end{itemize}

\begin{tcolorbox}[breakable,colback=teal!4,colframe=teal!35!black,title={Before you move on}]
\textquotedblleft{}In my opinion\textquotedblright{}? (\textbf{\sk{मम} \sk{मतेन}}.) \textquotedblleft{}I think that it will rain\textquotedblright{}? (\textbf{\sk{अहं} \sk{मन्ये} \sk{यत्} \sk{वर्षा} \sk{भविष्यति}}.) Where does \textbf{\sk{इति}} stand? (After the thought.)
\end{tcolorbox}
\section[\texorpdfstring{\sk{यतः} · \sk{कारणात्} (yataḥ · kāraṇāt)}{यतः · कारणात् (yataḥ · kāraṇāt)}]{\sk{यतः} · \sk{कारणात्} (yataḥ · kāraṇāt) — because, and because of}

\label{lesson:SA-C278-yatah-karanat}



\begin{quote}
Say \textquotedblleft{}I think that it will rain\textquotedblright{}. (\emph{ahaṃ manye yat varṣā bhaviṣyati}.) Now give the reason for a choice.
\end{quote}

\subsection*{You'll want to know: \sk{यतः}}
\textbf{\sk{यतः}} (\emph{yataḥ}), \textquotedblleft{}because\textquotedblright{}, opens the reason, and the reason can follow what you did:

\begin{quote}
\textbf{\sk{अहं} \sk{न} \sk{आगमिष्यामि},}
\end{quote}

\begin{quote}
\emph{ahaṃ na āgamiṣyāmi,}
\end{quote}

\begin{quote}
\textbf{\sk{यतः} \sk{अहं} \sk{रुग्णः} \sk{अस्मि।}}
\end{quote}

\begin{quote}
\emph{yataḥ ahaṃ rugṇaḥ asmi.} — \textquotedblleft{}I will not come, because I am ill.\textquotedblright{}
\end{quote}

\begin{quote}
\textbf{\sk{सः} \sk{न} \sk{पठति},}
\end{quote}

\begin{quote}
\emph{saḥ na paṭhati,}
\end{quote}

\begin{quote}
\textbf{\sk{यतः} \sk{सः} \sk{श्रान्तः} \sk{अस्ति।}}
\end{quote}

\begin{quote}
\emph{yataḥ saḥ śrāntaḥ asti.} — \textquotedblleft{}He is not reading, because he is tired.\textquotedblright{}
\end{quote}

\textbf{\sk{आगमिष्यामि}} (\emph{āgamiṣyāmi}) is \textbf{\sk{गमिष्यामि}} with the prefix \textbf{\sk{आ}-} of \textbf{\sk{आगच्छति}}: \textquotedblleft{}I will come\textquotedblright{}.

\subsection*{You'll want to know: \sk{कारणात्}, \sk{हेतुः}}
\textbf{\sk{कारणात्}} (\emph{kāraṇāt}) is \textbf{\sk{कारणम्}} (\emph{kāraṇam}), \textquotedblleft{}a reason\textquotedblright{}, with \emph{-āt}, \textquotedblleft{}from\textquotedblright{}: \textquotedblleft{}because of\textquotedblright{}. The cause goes in front of it, joined to it as one word:

\begin{quote}
\textbf{\sk{वर्षाकारणात्} \sk{सः} \sk{न} \sk{आगतवान्।}}
\end{quote}

\begin{quote}
\emph{varṣākāraṇāt saḥ na āgatavān.} — \textquotedblleft{}Because of the rain he did not come.\textquotedblright{}
\end{quote}

\textbf{\sk{हेतुः}} (\emph{hetuḥ}), \textquotedblleft{}a cause\textquotedblright{}, names the reason as a thing:

\begin{quote}
\textbf{\sk{हेतुः} \sk{किम्}?}
\end{quote}

\begin{quote}
\emph{hetuḥ kim?} — \textquotedblleft{}What is the cause?\textquotedblright{}
\end{quote}

\begin{quote}
\textbf{\sk{हेतुः} \sk{वर्षा} \sk{अस्ति।}}
\end{quote}

\begin{quote}
\emph{hetuḥ varṣā asti.} — \textquotedblleft{}The cause is the rain.\textquotedblright{}
\end{quote}

\subsection*{Your turn}
Say these aloud:

\begin{itemize}
\raggedright
  \item \textquotedblleft{}I will not come, because I am tired\textquotedblright{} — \emph{ahaṃ na āgamiṣyāmi, yataḥ ahaṃ śrāntaḥ asmi}
  \item \textquotedblleft{}I will not come, because I am ill\textquotedblright{} — \emph{ahaṃ na āgamiṣyāmi, yataḥ ahaṃ rugṇaḥ asmi}
  \item \textquotedblleft{}I will not come, because I am afraid\textquotedblright{} — \emph{ahaṃ na āgamiṣyāmi, yataḥ ahaṃ bhītaḥ asmi}
\end{itemize}

\begin{tcolorbox}[breakable,colback=teal!4,colframe=teal!35!black,title={Before you move on}]
\textquotedblleft{}Because\textquotedblright{} before a reason? (\textbf{\sk{यतः}}.) \textquotedblleft{}Because of the rain\textquotedblright{}? (\textbf{\sk{वर्षाकारणात्}}.) \textquotedblleft{}What is the cause?\textquotedblright{} (\textbf{\sk{हेतुः} \sk{किम्}?})
\end{tcolorbox}
\section[\texorpdfstring{\sk{मन्यते} (manyate)}{मन्यते (manyate)}]{\sk{मन्यते} (manyate) — thinks}

\label{lesson:SA-C278-manyate}



\begin{quote}
Before the new one: say \textquotedblleft{}I think that it will rain\textquotedblright{}. (\emph{ahaṃ manye yat varṣā bhaviṣyati}.) You have said this verb's \textquotedblleft{}I\textquotedblright{} form.
\end{quote}

\subsection*{You'll want to know: \sk{मन्यते}}
\textbf{\sk{मन्यते}} — \emph{manyate} — \textquotedblleft{}thinks\textquotedblright{}. \textbf{\sk{मन्ये}} (\emph{manye}) is its \textquotedblleft{}I\textquotedblright{}.

\begin{quote}
\textbf{\sk{सः} \sk{किं} \sk{मन्यते}?}
\end{quote}

\begin{quote}
\emph{saḥ kiṃ manyate?} — \textquotedblleft{}What does he think?\textquotedblright{}
\end{quote}

\begin{quote}
\textbf{\sk{मीरा} \sk{मन्यते} \sk{यत्} \sk{एतत्} \sk{सत्यम्।}}
\end{quote}

\begin{quote}
\emph{mīrā manyate yat etat satyam.} — \textquotedblleft{}Mira thinks that this is true.\textquotedblright{}
\end{quote}

\begin{cousinweb}
\textbf{\sk{मन्यते}} (\emph{manyate}) is built on the root √man, \textquotedblleft{}to think\textquotedblright{}, which Monier-Williams lists in the fourth class with the middle ending \emph{-te}. \textbf{\sk{मतम्}} (\emph{matam}), \textquotedblleft{}an opinion\textquotedblright{}, is its \emph{-ta} form: what is thought. It lives on as Hindi \emph{mānnā}, \textquotedblleft{}to respect, obey\textquotedblright{} (CDIAL 9857).
\end{cousinweb}

\subsection*{Your turn}
Say these aloud:

\begin{itemize}
\raggedright
  \item \emph{manyate}
  \item \textquotedblleft{}What does Mira think?\textquotedblright{} — \emph{mīrā kiṃ manyate?}
  \item \textquotedblleft{}Arun thinks that this is true\textquotedblright{} — \emph{aruṇaḥ manyate yat etat satyam}
\end{itemize}

\begin{tcolorbox}[breakable,colback=teal!4,colframe=teal!35!black,title={Before you move on}]
What does \textbf{\sk{मन्यते}} mean? (\textquotedblleft{}Thinks\textquotedblright{}.) Its \textquotedblleft{}I\textquotedblright{} form? (\textbf{\sk{मन्ये}}.)
\end{tcolorbox}
\section[\texorpdfstring{\sk{दृढः} (dṛḍhaḥ)}{दृढः (dṛḍhaḥ)}]{\sk{दृढः} (dṛḍhaḥ) — firm}

\label{lesson:SA-C278-drdhah}



\begin{quote}
Before the new one: say \textquotedblleft{}what does he think?\textquotedblright{} (\emph{saḥ kiṃ manyate?})
\end{quote}

\subsection*{You'll want to know: \sk{दृढः}}
\textbf{\sk{दृढः}} — \emph{dṛḍhaḥ} — \textquotedblleft{}firm\textquotedblright{}. It agrees with its noun: \textbf{\sk{दृढः}} for a masculine noun, \textbf{\sk{दृढा}} (\emph{dṛḍhā}) for a feminine one, \textbf{\sk{दृढम्}} (\emph{dṛḍham}) for a neuter one.

\begin{quote}
\textbf{\sk{मम} \sk{दृढः} \sk{विश्वासः} \sk{अस्ति।}}
\end{quote}

\begin{quote}
\emph{mama dṛḍhaḥ viśvāsaḥ asti.} — \textquotedblleft{}I have a firm belief.\textquotedblright{}
\end{quote}

\begin{quote}
\textbf{\sk{गृहं} \sk{दृढम्} \sk{अस्ति।}}
\end{quote}

\begin{quote}
\emph{gṛhaṃ dṛḍham asti.} — \textquotedblleft{}The house is solid.\textquotedblright{}
\end{quote}

\begin{quote}
\textbf{\sk{मम} \sk{निर्णयः} \sk{दृढः।}}
\end{quote}

\begin{quote}
\emph{mama nirṇayaḥ dṛḍhaḥ.} — \textquotedblleft{}My decision is firm.\textquotedblright{}
\end{quote}

\begin{cousinweb}
\textbf{\sk{दृढः}} (\emph{dṛḍhaḥ}) comes from \textbf{\sk{दृढ}} (\emph{dṛḍha}), \textquotedblleft{}fixed, firm, strong\textquotedblright{}, in Monier-Williams; CDIAL 6508 sets it under the root √dṛṃh and lists Hindi \emph{diṛh}, \textquotedblleft{}firm\textquotedblright{}, among its descendants.
\end{cousinweb}

\subsection*{Your turn}
Say these aloud:

\begin{itemize}
\raggedright
  \item \emph{dṛḍhaḥ}
  \item \textquotedblleft{}My decision is firm\textquotedblright{} — \emph{mama nirṇayaḥ dṛḍhaḥ}
  \item \textquotedblleft{}In my opinion the house is solid\textquotedblright{} — \emph{mama matena gṛhaṃ dṛḍham}
\end{itemize}

\begin{tcolorbox}[breakable,colback=teal!4,colframe=teal!35!black,title={Before you move on}]
What does \textbf{\sk{दृढः}} mean? (\textquotedblleft{}Firm\textquotedblright{}.) \textquotedblleft{}A solid house\textquotedblright{}? (\textbf{\sk{दृढं} \sk{गृहम्}}.)
\end{tcolorbox}
\section[\texorpdfstring{\sk{ढ} (ḍha)}{ढ (ḍha)}]{\sk{ढ} — one character, met first in \sk{दृढः}}

\label{lesson:SA-S227-letter-ddha}



\begin{quote}
Before the new one: say \textbf{\sk{दृढः}} (\emph{dṛḍhaḥ}), \textquotedblleft{}firm\textquotedblright{}. Then say the sounds of \textbf{\sk{ड}} and \textbf{\sk{घ}}. (\emph{ḍa}, \emph{gha}.)

One character this time: the one in the middle of \textbf{\sk{दृढः}}.
\end{quote}

\subsection*{Script you'll notice: \sk{ढ}}
\textbf{\sk{ढ}} — \emph{ḍha}.

It is a \textbf{consonant}: \textbf{\sk{ड}} (\emph{ḍa}) with a breath after it, as \textbf{\sk{घ}} is \textbf{\sk{ग}} with a breath.

What it is made of:

\begin{itemize}
\raggedright
  \item a descending stem joined to a broad outer bowl and closed inner loop
  \item the top shirorekhā
\end{itemize}

You have met it here:

\begin{itemize}
\raggedright
  \item \textbf{\sk{दृढः}} \emph{dṛḍhaḥ} — firm
\end{itemize}

\subsection*{Writing: \sk{ढ}}
\hlblockfigure{\detokenize{figures/SA-S227-letter-ddha-filmstrip.pdf}}{How \sk{ढ} is written, stroke by stroke}

\begin{itemize}
\raggedright
  \item \textbf{1.} start at the stem's headline junction, descend, turn left through the upper shoulder, sweep counterclockwise around the broad outer bowl, and continue counterclockwise around the closed inner loop without lifting
  \item \textbf{2.} lift and draw the top shirorekhā left-to-right
\end{itemize}

\textbf{Pen lifts: 1.} The pen comes up 1 time and no more.

\begin{quote}
Verified two-stroke teaching form fitted to the bundled printed outline.
\end{quote}

\begin{quote}
This is one attested teaching order and not a national standard — handwriting here is taught with school-to-school variation. Source: Opiaterein, ‘Deva-\sk{ढ}-order.gif’, strokes 1–2, Wikimedia Commons, 10 May 2009.
\end{quote}

\subsection*{Your turn}
\begin{itemize}
\raggedright
  \item \emph{Look:} at this, and find \sk{ढ} in it
\end{itemize}

\begin{quote}
\sk{दृढः}
\end{quote}

\begin{itemize}
\raggedright
  \item \emph{Trace it:} \sk{ढ} three times, saying \emph{ḍha} as you finish each one
  \item \emph{Look:} at \textbf{\sk{मम} \sk{निर्णयः} \sk{दृढः}} and find \sk{ढ} once more
\end{itemize}

\begin{tcolorbox}[breakable,colback=teal!4,colframe=teal!35!black,title={Before you move on}]
Which character is this — \sk{ढ}? What sound does it carry? (\emph{ḍha}.) Which word did it come from? (\textbf{\sk{दृढः}}, \textquotedblleft{}firm\textquotedblright{}.)
\end{tcolorbox}
\section[(review)]{Review — what I think, and why}

\label{lesson:SA-R278-opinion-and-reason}



\begin{quote}
Say \textquotedblleft{}first Mira went to the shop\textquotedblright{}. (\emph{prathamaṃ mīrā āpaṇaṃ gatavatī}.) Then \textquotedblleft{}she came in and sat down\textquotedblright{}. (\emph{sā praviśya upaviśati}.)
\end{quote}

\subsection*{Your turn}
\begin{itemize}
\raggedright
  \item \emph{Say it:} In my opinion the city is beautiful.
  \item \emph{Say it:} I think that Mira will come tomorrow.
  \item \emph{Say it:} I will not go, because I am tired.
  \item \emph{Say it:} Because of the rain Mira did not come.
  \item \emph{Say it:} What does Arun think? My decision is firm.
  \item \emph{Read it:} \textbf{\sk{दृढः}} and \textbf{\sk{मम} \sk{दृढः} \sk{विश्वासः} \sk{अस्ति}} aloud
\end{itemize}

\begin{tcolorbox}[breakable,colback=teal!4,colframe=teal!35!black,title={Before you move on}]
\textbf{\sk{मन्ये}} or \textbf{\sk{मन्यते}}, for \textquotedblleft{}I think\textquotedblright{}? (\textbf{\sk{मन्ये}}.) \textbf{\sk{यतः}} or \textbf{\sk{कारणात्}}, before a whole reason? (\textbf{\sk{यतः}}.) Now the other way round: what does \textbf{\sk{हेतुः} \sk{किम्}?} ask? (\textquotedblleft{}What is the cause?\textquotedblright{})
\end{tcolorbox}
